% a06_BALtorus_A4: continuation of a06_BALtorus_A3 (input from it). Group (A), part 4 of 4: the torus.

\begin{lemma}[{\cite[Prop.~2.1]{Ba1} on a torus, with torus-free constants}]\label{lem:a06A:torus}
Let $T$ be a torus $\eta\Z^d/\Gamma$ whose period lattice $\Gamma$ is contained in $M\Z^d$ (so the large blocks of the unit lattice tile $T$) and whose sup-systole is at least $3M$; let $\Omega\subset T$ be a union of large blocks ($\Omega=T$ included), $A$ regular on $\Omega$, and $e_k$ small. Then the theorem of \cite{Ba4} (\cite[Prop.~2.1, (2.24)--(2.26)]{Ba1}) holds on $T$ in its restricted form, with distances the quotient sup-distances and with the constants $\delta_0$, $c_0$, $R_0$ of \cite{Ba4}, which do not depend on $\Gamma$; for $\Omega=T$ it holds for all $x,x'\in T$. The same is true for tori of sup-systole less than $3M$ (but at least $M$) whose period lattice lies in $M\Z^d$. Consequently every estimate of \cref{thm:a06A:groupA} holds on $T$ with the constants of \cref{thm:a06A:groupA}.
\end{lemma}

\begin{proof}
The proof of \cite[Sect.~2]{Ba4} (pp.~574--580) is carried onto $T$ step by step; it uses the geometry of the lattice in the following places, and nowhere else.

(1) \emph{The local cubes.} Let $J$ index the vertices $Mj$ of the large-block partition of $T$ and $C_j\defeq Mj+[-M,M]^d$ (in unit-lattice units). Since $2M$ is less than the sup-systole, the projection $\eta\Z^d\to T$ is injective on each $C_j$; it preserves bonds, unit blocks, $L$-blocks and the contours of $Q_k$ inside blocks. So $\Pi_j\defeq\Omega\cap C_j$ \cite[p.~575]{Ba4} is a union of at most $2^d$ large blocks inside an embedded cube of side $2M$, a cube when $Mj$ is not a boundary point of $\Omega$, and the decomposition $A=A_0+A'$ on $C_j$ with the constant $A_0\defeq A(Mj)$ and the local field $A_j\defeq A_0+\theta_jA'$ of \cite[p.~575]{Ba4} is built from $A|_{C_j}$ alone, with $|A'|,|\partial^\eta A'|\le c'e_k^{\beta-1}$, $c'=dMc$. Hence \cite[Lemma~2.1]{Ba4} applies to every $G_k(\Pi_j,A_j)$ and \cite[Lemma~2.2]{Ba4} to the interior cubes (for $\Pi_j$ meeting $\partial\Omega$, $A_j\defeq A$ \cite[p.~575]{Ba4}), with their constants, which depend on $d$, $M$, $\alpha$, $c$ only; \cite[Lemma~2.4]{Ba4}, used in the proof of \cite[Lemma~2.2]{Ba4}, is \cref{thm:a06:A5} on parallelepipeds of the fine lattice.

(2) \emph{The partition of unity.} With $h\in C_0^\infty((-\frac23,\frac23))$, $h=1$ on $[-\frac13,\frac13]$, $\sum_{n\in\Z}h(\cdot-n)^2=1$, and $h_j(x)\defeq\prod_\mu h(x_\mu/M-j_\mu)$ \cite[p.~575]{Ba4}: because $\Gamma\subset M\Z^d$ the family $\{h_j\}_{j\in J}$ is well defined on $T$, $\sum_{j\in J}h_j^2=1$ on $T$ and $\supp h_j\subset C_j$. So the approximate inverse $G_0=\sum_jh_jG_k(\Pi_j,A_j)h_j$ \cite[(2.2)]{Ba4} and the identities \cite[(2.3)--(2.9)]{Ba4} hold on $T$ as on $\eta\Z^d$; for $\Omega=T$ the boundary of $\Omega$ is empty and the boundary term of \cite[(2.6)]{Ba4} vanishes.

(3) \emph{The remainder.} Each $j$ has at most $3^d-1$ neighbours and each point lies in at most $2^d$ supports, so the operator $R$ of \cite[(2.11)]{Ba4} has $\norm R_{2,2}\le c_dO(M^{-1})<1$ by \cite[Lemma~2.1]{Ba4}, as on \cite[p.~577]{Ba4}; the Neumann series \cite[(2.12)]{Ba4} converges and gives the walk representation \cite[(2.13)]{Ba4}, the walks being sequences of indices $j\in J$ with consecutive terms adjacent, and the count $2^d(3^d)^{n-1}2^d$ of walks of length $n$ \cite[p.~579]{Ba4} is unchanged.

(4) \emph{The metric.} Consecutive cubes of a walk have centres at quotient distance at most $M$, so a walk of length $n$ from a cube containing $x,x'$ to a cube containing points of $\supp f$ has $\dist(\{x,x'\},\supp f)\le(n+2)M$, i.e.\ $n\ge M^{-1}\dist(\{x,x'\},\supp f)-2$, which is the restriction used in \cite[(2.22)]{Ba4}; with (3), the geometric sum \cite[(2.22)]{Ba4} gives $\delta_0$ and $c_0$ as on $\eta\Z^d$, once $M$ is fixed with $3^dc_{2.1}O(1)M^{-1}\le\mathrm e^{-1}$, $c_{2.1}$ the constant of \cite[Lemma~2.1, (2.15)]{Ba4} (not the $c_2$ of \cref{lem:a06A:blockinv}). The definition \cite[(2.19)]{Ba4} of $R_0$ (the diameter of the union of the first $n_0$ cubes of a walk, plus $2M$) and the implication ``$\dist(\{x,x'\},\Omega^c)\ge R_0$, hence the first $n_0$ cubes $\Pi_{\omega_i}$ of every walk from $x,x'$ are cubes contained in $\Omega$'' use only the quotient metric and (1); on these cubes $A_j$ is constant near $\partial\Pi_j$ and \cite[Lemma~2.2]{Ba4} applies, on the others only \cite[Lemma~2.1]{Ba4} is used, as in \cite[(2.20)--(2.21)]{Ba4}. The cancellation of the terms with interior cubes in the two representations of $\delta G_k(\Omega,\Omega_0,A)$ \cite[p.~579]{Ba4} uses (1)--(2) only. If $\Omega=T$, every $\Pi_j=C_j$ is a cube with $A_j$ constant near its boundary, so no restriction on $x,x'$ is needed \cite[p.~611]{Ba1}.

(5) \emph{Small tori.} If the sup-systole of $T=\eta\Z^d/\Gamma$ is less than $3M$, let $\tilde T\defeq\eta\Z^d/3\Gamma$, a covering torus of systole at least $3M$ with deck group $(\Z/3)^d$. The operators $-\Delta^{\eta,N}_{A,\Omega}+m_k^2+a_kP_k(A)$ lifted to $\tilde T$ commute with the deck translations and map deck-invariant functions to deck-invariant functions, so for $f$ on $\Omega$ the lift of $G_k(\Omega,A)f$ is $G_k(\tilde\Omega,\tilde A)\tilde f$, with $\tilde\Omega,\tilde A,\tilde f$ the lifts; $\norm{\tilde f}_\infty=\norm f_\infty$, and the quotient distances in $\tilde T$ from a lift of $x$ to $\supp\tilde f$ and to $\tilde\Omega^c$ equal those in $T$. So the bounds on $\tilde T$, proved in (1)--(4), descend to $T$ with the same constants. The same applies to $\delta G_k$.

No other constant of \cite[Sect.~2]{Ba4} involves the global geometry. Finally, \cref{lem:a06A:cube,lem:a06A:coer,lem:a06A:CT,lem:a06A:blockinv,lem:a06A:resdiff,lem:a06A:CLdiff}, the first bound of \cref{prop:a06A:Ba1P22} and \cref{lem:a06A:polylog} are proved above on an arbitrary union of blocks of an arbitrary torus (their only global ingredients are the sums $S(\theta)$ and $S_L(\theta)$ over the block lattice, bounded by their values on $\Z^d$), and \cref{lem:a06A:twosided,lem:a06A:domdiff} and the second bound of \cref{prop:a06A:Ba1P22} use the printed theorem only in the form just established. Hence \cref{thm:a06A:groupA} holds on $T$ with the same constants.
\end{proof}
